% ============================================================
% HyperCut 项目汇报 Beamer 幻灯片
% 编译方式：XeLaTeX（支持中文 ctex）
%   本地：xelatex presentation.tex
%   Overleaf：上传后选择 XeLaTeX 编译器
% 需要 tikz / listings 宏包（TeX Live / Overleaf 自带）
% ============================================================
\documentclass[aspectratio=169,11pt]{beamer}

\usepackage[UTF8]{ctex}
\usepackage{tikz}
\usetikzlibrary{positioning, arrows.meta, shapes.geometric, calc}
\usepackage{listings}

% ---------- 主题与配色（青色调，呼应前端 UI） ----------
\usetheme{Madrid}
\usecolortheme{default}
\definecolor{teal}{RGB}{40,127,131}
\definecolor{tealDark}{RGB}{22,94,98}
\definecolor{lime}{RGB}{200,242,92}
\definecolor{ink}{RGB}{32,36,40}
\setbeamercolor{structure}{fg=teal}
\setbeamercolor{title}{fg=white, bg=teal}
\setbeamercolor{frametitle}{fg=white, bg=teal}
\setbeamercolor{block title}{fg=white, bg=teal}
\setbeamercolor{block body}{bg=teal!8}
\setbeamercolor{item}{fg=tealDark}
\setbeamercolor{section in toc}{fg=ink}
\setbeamertemplate{navigation symbols}{}
\setbeamertemplate{footline}[frame number]

\lstset{basicstyle=\ttfamily\scriptsize, keywordstyle=\color{tealDark}\bfseries,
  commentstyle=\color{gray}, stringstyle=\color{teal}, frame=single, rulecolor=\color{teal!40},
  breaklines=true, columns=fullflexible, showstringspaces=false}

\title{代码驱动视频剪辑 Agent\\—— HyperCut 系统设计与实现}
\author{汇报人：Jiahao Liu}
\institute{项目：HyperCut · π Agent $\times$ HyperFrames}
\date{\today}

\begin{document}

% ============ 封面 ============
\begin{frame}
  \titlepage
\end{frame}

% ============ 目录 ============
\begin{frame}{目录}
  \tableofcontents
\end{frame}

% ============ 第 1 部分：项目概述 ============
\section{项目概述}

\begin{frame}{项目要解决的问题}
  \begin{block}{核心目标}
    让用户用\textbf{一句话}描述需求，系统自动生成并渲染出一段\textbf{MP4 视频}。
  \end{block}
  \vspace{0.5em}
  传统视频剪辑痛点：
  \begin{itemize}
    \item 需要在时间线上手动拖拽素材、逐段剪辑
    \item 门槛高、耗时长，重复性工作多
  \end{itemize}
  \vspace{0.5em}
  我们的思路——\textbf{代码驱动剪辑}：
  \begin{itemize}
    \item 用代码（HTML）描述视频：何时出现什么画面、动画怎么做、字幕如何同步
    \item 由 Agent 自动生成代码，渲染引擎逐帧输出视频
    \item 代码是确定性的：同样输入永远得到同样成片
  \end{itemize}
\end{frame}

\begin{frame}{三大设计主线}
  \begin{columns}[T]
    \begin{column}{0.32\textwidth}
      \begin{block}{① 代码驱动剪辑}
        用 HyperFrames HTML 描述时间线，逐帧渲染成片，可复现、可版本化。
      \end{block}
    \end{column}
    \begin{column}{0.32\textwidth}
      \begin{block}{② Agent 驱动}
        π Agent 理解意图、调用工具、生成代码，形成「思考—执行—观察」闭环。
      \end{block}
    \end{column}
    \begin{column}{0.32\textwidth}
      \begin{block}{③ 人机交互}
        Web 前端：输入需求、选择风格、实时预览、增量修改、音频合成。
      \end{block}
    \end{column}
  \end{columns}
\end{frame}

% ============ 第 2 部分：总体架构 ============
\section{总体架构}

\begin{frame}{系统三层架构}
  \begin{center}
    \begin{tikzpicture}[
        box/.style={draw=teal, fill=teal!8, rounded corners, align=center, minimum width=3.6cm, minimum height=1.3cm, font=\small},
        arr/.style={-{Stealth[length=3mm]}, thick, tealDark}
      ]
      \node[box, fill=teal!20] (fe) {前端 Web 应用\\React 19 + Vite + TS};
      \node[box, below=0.9cm of fe, fill=lime!25] (be) {后端服务\\π Agent + Node.js};
      \node[box, below=0.9cm of be, fill=teal!8] (ren) {渲染引擎\\HyperFrames + FFmpeg};
      \draw[arr] (fe) -- node[right, font=\scriptsize]{REST API} (be);
      \draw[arr] (be) -- node[right, font=\scriptsize]{生成 HTML / 调用 CLI} (ren);
      \draw[arr] ([xshift=-0.6cm]be.west) -- node[left, font=\scriptsize]{轮询状态} ([xshift=-0.6cm]fe.west);
    \end{tikzpicture}
  \end{center}
  \vspace{0.3em}
  \begin{itemize}
    \item \textbf{前端}：用户交互层，输入需求、选择风格/时长、预览成片、编辑音频
    \item \textbf{后端}：π Agent 运行时 + 业务逻辑，连接 LLM 与渲染引擎
    \item \textbf{渲染}：HyperFrames 把 HTML 合成页逐帧渲染成 MP4
  \end{itemize}
\end{frame}

\begin{frame}{技术栈}
  \begin{columns}[T]
    \begin{column}{0.5\textwidth}
      \textbf{前端}
      \begin{itemize}
        \item React 19 + TypeScript
        \item Vite 8（构建 / 热更新）
        \item lucide-react（图标）
      \end{itemize}
      \vspace{0.8em}
      \textbf{后端}
      \begin{itemize}
        \item Node.js（原生 http 服务）
        \item π Agent：\texttt{pi-agent-core} + \texttt{pi-ai}
        \item OpenAI 兼容端点（gpt-5.6-sol）
      \end{itemize}
    \end{column}
    \begin{column}{0.5\textwidth}
      \textbf{渲染与媒体}
      \begin{itemize}
        \item HyperFrames 0.8.x（HTML → 视频）
        \item FFmpeg / FFprobe（编码与音频分析）
      \end{itemize}
      \vspace{0.8em}
      \textbf{关键特性}
      \begin{itemize}
        \item 代码驱动、确定性渲染
        \item 支持中文提示词
        \item 风格模板、分段长视频、音频轨
      \end{itemize}
    \end{column}
  \end{columns}
\end{frame}

% ============ 第 3 部分：后端实现 ============
\section{后端实现}

\begin{frame}{π Agent 运行时}
  \begin{columns}[T]
    \begin{column}{0.5\textwidth}
      \textbf{Agent 循环}（核心思想）
      \begin{center}
        \begin{tikzpicture}[
            n/.style={draw=teal, fill=teal!10, rounded corners, minimum width=2.4cm, minimum height=0.7cm, align=center, font=\scriptsize},
            arr/.style={-{Stealth[length=2mm]}, tealDark, thick}
          ]
          \node[n] (t) {一句任务};
          \node[n, below=0.5cm of t] (m) {模型思考};
          \node[n, below=0.5cm of m] (tool) {调用工具};
          \node[n, below=0.5cm of tool] (o) {观察结果};
          \draw[arr] (t) -- (m);
          \draw[arr] (m) -- (tool);
          \draw[arr] (tool) -- (o);
          \draw[arr] (o.east) to[out=0,in=0] node[right, font=\scriptsize]{没完成？带结果回「思考」} (m.east);
        \end{tikzpicture}
      \end{center}
    \end{column}
    \begin{column}{0.5\textwidth}
      \textbf{我们的工具集}
      \begin{itemize}
        \item \texttt{write\_composition} —— 写/改合成页 HTML
        \item \texttt{modify\_composition} —— 帧级增量修改
        \item 风格注入、分段上下文
      \end{itemize}
      \vspace{0.4em}
      \textbf{关键设计}
      \begin{itemize}
        \item SYSTEM\_PROMPT 定义 HyperFrames 硬约束
        \item 工具参数用 TypeBox schema 校验
        \item 超时 + 自动终止（5 分钟）
      \end{itemize}
    \end{column}
  \end{columns}
\end{frame}

\begin{frame}[fragile]{核心：HyperFrames 合成页（代码即时间线）}
  \begin{columns}[T]
    \begin{column}{0.5\textwidth}
      HTML 用 \texttt{data-*} 属性声明时序：
      \lstset{language=HTML, basicstyle=\ttfamily\tiny}
      \begin{lstlisting}
<div id="stage" class="stage"
  data-composition-id="hypercut-main"
  data-no-timeline data-start="0"
  data-duration="30" data-width="960"
  data-height="540" data-fps="24">
  <h1 class="clip" data-start="0"
      data-duration="6" data-track-index="1">
    标题文字
  </h1>
</div>
      \end{lstlisting}
    \end{column}
    \begin{column}{0.5\textwidth}
      动画用 \texttt{hf-seek} 事件驱动（确定性）：
      \lstset{language=JavaScript, basicstyle=\ttfamily\tiny}
      \begin{lstlisting}
window.addEventListener('hf-seek',
  event => {
    const t = event.detail.time;
    // 由 t 计算画面状态，无随机、无定时器
    title.style.opacity =
      Math.min(1, t / 1.2);
    // 确定性：同一 t 永远同一帧
  });
      \end{lstlisting}
      \vspace{0.4em}
      \begin{block}{为什么用代码？}
        确定性、可版本化、天然适合让 Agent 生成与修改。
      \end{block}
    \end{column}
  \end{columns}
\end{frame}

\begin{frame}{后端模块划分}
  \begin{center}
    \begin{tikzpicture}[
        f/.style={draw=teal, fill=teal!8, rounded corners, align=center, minimum width=3.2cm, minimum height=0.85cm, font=\scriptsize},
        arr/.style={-{Stealth[length=2mm]}, tealDark}
      ]
      \node[f] (index) {\texttt{server/index.mjs}\\HTTP 服务 + 路由 + 任务管理};
      \node[f, below=0.7cm of index] (agent) {\texttt{pi-agent.mjs}\\π Agent 封装 + 提示词};
      \node[f, below=0.7cm of agent] (style) {\texttt{styles.mjs}\\风格模板 + 意图识别};
      \node[f, below=0.7cm of style] (audio) {\texttt{audio-store.mjs}\\音频素材与音轨};
      \draw[arr] (index) -- (agent);
      \draw[arr] (agent) -- (style);
      \draw[arr] (style) -- (audio);
    \end{tikzpicture}
  \end{center}
  \vspace{0.4em}
  \begin{itemize}
    \item \texttt{index.mjs}：原生 http 服务，处理 REST API、任务队列、渲染调度
    \item \texttt{pi-agent.mjs}：创建 π Agent，注入 SYSTEM\_PROMPT 与工具
    \item \texttt{styles.mjs}：8 套风格模板（palette/keywords/详细 prompt），纯数据、可模块化替换
    \item \texttt{audio-store.mjs}：音频导入、ffprobe 分析、音轨参数管理
  \end{itemize}
\end{frame}

% ============ 第 4 部分：核心流程与功能 ============
\section{核心流程与功能}

\begin{frame}{从一句话到成片：完整流水线}
  \begin{center}
    \begin{tikzpicture}[
        n/.style={draw=teal, fill=teal!10, rounded corners, minimum width=2.1cm, minimum height=0.9cm, align=center, font=\scriptsize},
        arr/.style={-{Stealth[length=2.5mm]}, tealDark, thick}
      ]
      \node[n] (a) {用户输入需求};
      \node[n, right=0.5cm of a] (b) {需求理解\\标题/风格/时长};
      \node[n, right=0.5cm of b] (c) {Agent 生成\\HyperFrames HTML};
      \node[n, right=0.5cm of c] (d) {逐帧渲染};
      \node[n, right=0.5cm of d] (e) {MP4 成片};
      \draw[arr] (a) -- (b);
      \draw[arr] (b) -- (c);
      \draw[arr] (c) -- (d);
      \draw[arr] (d) -- (e);
    \end{tikzpicture}
  \end{center}
  \vspace{0.5em}
  \begin{enumerate}
    \item \textbf{需求理解}：正则 + 关键词解析标题、配色、总时长（支持「一分钟/60秒/1分30秒」）
    \item \textbf{生成代码}：π Agent 依据风格模板与分段信息生成完整 HTML
    \item \textbf{逐帧渲染}：HyperFrames 用 headless Chrome 截图 + FFmpeg 编码
    \item \textbf{成片}：MP4 写盘，前端播放器展示
  \end{enumerate}
\end{frame}

\begin{frame}[fragile]{功能①：风格模板系统（意图自动识别）}
  \begin{columns}[T]
    \begin{column}{0.55\textwidth}
      \textbf{8 套风格模板}，每套含三要素：
      \begin{itemize}
        \item \texttt{palette}：主色/辅色/底色
        \item \texttt{keywords}：用于意图自动识别
        \item \texttt{template}：详细的七段式 prompt（配色/排版/构图/背景/动效/元素/避免）
      \end{itemize}
    \end{column}
    \begin{column}{0.45\textwidth}
      \textbf{自动识别（关键词路由）}：
      \begin{lstlisting}[language=JavaScript, basicstyle=\ttfamily\tiny]
function detectStyle(prompt){
  let best=null, score=0;
  for(const s of STYLES){
    let n=0;
    for(const kw of s.keywords)
      if(prompt.includes(kw)) n++;
    if(n>score){ best=s; score=n; }
  }
  return best ?? STYLES[0];
}
      \end{lstlisting}
    \end{column}
  \end{columns}
  \vspace{0.4em}
  \begin{exampleblock}{示例}
    「做一个数学公式讲解视频并剪辑」 $\rightarrow$ 自动识别为「数学学术」风格，零额外 LLM 调用。
  \end{exampleblock}
\end{frame}

\begin{frame}{功能②：分段生成长视频（先预览后拼接）}
  \begin{center}
    \begin{tikzpicture}[
        n/.style={draw=teal, fill=teal!10, rounded corners, minimum width=2.6cm, minimum height=0.9cm, align=center, font=\scriptsize},
        arr/.style={-{Stealth[length=2.5mm]}, tealDark, thick}
      ]
      \node[n] (total) {60 秒需求};
      \node[n, below=0.6cm of total] (s1) {段1：0--30s\\生成并渲染};
      \node[n, right=0.5cm of s1] (preview) {首段立即预览\\(缩短等待)};
      \node[n, below=0.6cm of s1] (s2) {段2：30--60s\\延续段1生成};
      \node[n, below=0.6cm of s2] (concat) {FFmpeg 拼接\\完整 60s};
      \draw[arr] (total) -- (s1);
      \draw[arr] (s1) -- (preview);
      \draw[arr] (s1) -- (s2);
      \draw[arr] (s2) -- (concat);
    \end{tikzpicture}
  \end{center}
  \begin{itemize}
    \item 按 30 秒切段，逐段「生成—渲染—预览」，首段完成即推给前端
    \item 全部段完成后用 FFmpeg \texttt{concat} 无损拼接
    \item \textbf{连贯性}：后段生成时，把前段完整 HTML 作为上下文传入，提示「延续上一段的配色/构图/叙事结尾」
  \end{itemize}
\end{frame}

\begin{frame}{功能③：音频轨 与 功能④：实时修改}
  \begin{columns}[T]
    \begin{column}{0.5\textwidth}
      \begin{block}{音频轨}
        \begin{itemize}
          \item 导入 MP3/WAV/M4A/OGG，ffprobe 自动分析时长/码率
          \item 音轨参数：音量、淡入淡出、循环、静音、素材偏移
          \item 以 \texttt{<audio>} 元素注入合成页，重新渲染
        \end{itemize}
      \end{block}
    \end{column}
    \begin{column}{0.5\textwidth}
      \begin{block}{实时修改（帧级增量编辑）}
        \begin{itemize}
          \item 指定某一帧，Agent 只改该帧之后的代码
          \item 保留无关内容，增量生成新版本
          \item 版本号递增，可追溯
        \end{itemize}
      \end{block}
    \end{column}
  \end{columns}
\end{frame}

% ============ 第 5 部分：前端实现 ============
\section{前端实现}

\begin{frame}[fragile]{前端架构（React 单页应用）}
  \begin{columns}[T]
    \begin{column}{0.55\textwidth}
      \textbf{界面布局}
      \begin{itemize}
        \item 左侧：项目面板（任务/版本/服务状态）
        \item 中间：预览画布 + 时间线
        \item 右侧：四栏 Inspector
      \end{itemize}
      \vspace{0.4em}
      \textbf{四个 Tab}
      \begin{enumerate}
        \item 生成设置：描述、风格、时长
        \item 音频：导入与音轨配置
        \item 代码：查看生成的 HTML
        \item 活动记录：操作历史
      \end{enumerate}
    \end{column}
    \begin{column}{0.45\textwidth}
      \textbf{核心状态}
      \begin{lstlisting}[language=JavaScript, basicstyle=\ttfamily\tiny]
const [prompt, setPrompt] = ...
const [duration, setDuration] = ... // 时长
const [styleId, setStyleId] = ... // 风格
const [videoUrl, setVideoUrl] = ... // 成片
const [phase, setPhase] = ... // 阶段
const [audioTrack, setAudioTrack] = ...
      \end{lstlisting}
    \end{column}
  \end{columns}
\end{frame}

\begin{frame}[fragile]{前后端交互（REST API + 轮询）}
  \begin{columns}[T]
    \begin{column}{0.5\textwidth}
      \textbf{接口清单}
      \begin{itemize}
        \item \texttt{POST /api/agent/generate} —— 生成
        \item \texttt{POST /api/agent/modify} —— 修改
        \item \texttt{GET /api/styles} —— 风格列表
        \item \texttt{GET/PUT/DELETE /api/audio-track}
        \item \texttt{POST /api/assets} —— 上传音频
        \item \texttt{GET /api/jobs/:id} —— 任务状态
        \item \texttt{GET /api/video/:id} —— 播放成片
      \end{itemize}
    \end{column}
    \begin{column}{0.5\textwidth}
      \textbf{状态轮询}：前端每 1 秒轮询任务状态，据此更新进度条与预览。
      \begin{lstlisting}[language=JavaScript, basicstyle=\ttfamily\tiny]
async function poll(jobId){
  for(let i=0;i<1800;i++){
    await sleep(1000);
    const job = await fetch(
      `/api/jobs/${jobId}`).then(r=>r.json());
    if(job.videoUrl) setVideoUrl(...);
    if(job.status === "complete") break;
    if(job.status === "failed") throw ...;
  }
}
      \end{lstlisting}
    \end{column}
  \end{columns}
\end{frame}

% ============ 第 6 部分：总结 ============
\section{总结与展望}

\begin{frame}{项目亮点总结}
  \begin{block}{已实现的核心能力}
    \begin{itemize}
      \item \textbf{代码驱动剪辑}：HyperFrames HTML 描述时间线，确定性逐帧渲染
      \item \textbf{Agent 驱动}：π Agent 循环 + 工具调用，自然语言 → 成片
      \item \textbf{风格模板}：8 套风格 + 意图自动识别，可模块化扩展
      \item \textbf{分段长视频}：先预览后拼接，缩短等待，前后段连贯
      \item \textbf{音频轨 + 实时修改}：BGM 合成、帧级增量编辑、版本管理
    \end{itemize}
  \end{block}
\end{frame}

\begin{frame}{不足与展望}
  \begin{columns}[T]
    \begin{column}{0.5\textwidth}
      \textbf{当前不足}
      \begin{itemize}
        \item 分段拼接依赖分辨率/帧率一致，失败时需重编码回退
        \item 意图识别为关键词规则，复杂语义理解有限
        \item 逐帧渲染耗时较长
      \end{itemize}
    \end{column}
    \begin{column}{0.5\textwidth}
      \textbf{后续展望}
      \begin{itemize}
        \item 引入语义级风格理解（LLM 分类/生成）
        \item 先出全片大纲，再逐段按大纲生成
        \item 支持真实视频素材导入与多模态理解
        \item 云渲染 / 分布式渲染加速
      \end{itemize}
    \end{column}
  \end{columns}
  \vspace{1em}
  \begin{center}
    \Large 谢谢！欢迎提问
  \end{center}
\end{frame}

\end{document}
